\subsection{PARTISANS}

Confederate Partisans (guerrillas) were particularly numerous and effective in the Shenandoah Valley region. Composed of frces numbering less than a dozen to forces of batallion size, the Partisans varied from troops who were regularly enrolled in the Confederate army to bands of brigands that no one really wished to claim.

They performed as scouts, sometimes fought in formal battles alongside regular units, constantly harassed the Union supply lines, and made innumerable raids on the B\&O Railroad. They tied up considerable numbers of Union garrison troops throughout the war. Their forte was the surprise attack out of nowhere, and they proved extremely difficult to catch (often adopting the guise of harmless farmers when pressed by Union pursuits).

These Shenandoah Partisans are represented in the game by three unit countesr, each counter containing the name of one of the most successful Partisan commanders.

\begin{enumerate}
  \item Partisans can largely ignore supply considerations. They are always considered to be supplied unless they are in a Devastated hex. In a Devastated hex they must be considered "unsupplied" unless within the normal four hex distance of a friendly Supply counter.

  \item Partisans can be removed completely from the Map when not adjacent to an enemy unit. The procedure for this is as follows:

  \begin{enumerate}[label=(\alph*)]
    \item The Confederate player moves the Partisans normally until over two hexes from any Union units. He then asks the Union Player to avert his head for a minute and not look at the Map as he completes his movement. The Confederate Player then finishes moving the Partisan(s), and writes on a sheet of paper the grid coordinates of the hex he is in at the end of movement.

    \item When complete, the sheet of paper is folded to conceal the writing and placed in some convenient location near the Union Player, who is not yet allowed to look at it.

    \item At the start of the Confederate Player's next Phase, the Union Player unfolds the note and places the Partisan unit where noted. The Confederate Player then again begins to move, and Steps (a) through (c) are repeated on every Turn.
  \end{enumerate}

  \item The Union Player must be told immediately if he enters the Zone of Control of an off-board Partisan unit while he is moving. In this case the Partisan unit is placed back on the Map, and normal Combat takes place (Note: if the Union units lack the movement to enter the Partisan Zone of Control, they must end their Movem in the last hex they were in prior to attempting to enter the hidden Zone of Control).

  \item Partisan units are treated in all other respects like normal Cavalry units.

  \item When a Partisan unit is eliminated, the Union Player gets Cavalry value Victory Points. However, the Partisan unit can be re-introduced to the Mapboard seven Turns after its elimination, again being set in any hex two hexes from the nearest Union unit.
\end{enumerate}